\subsubsection{An optimal two-moment contour shield}
\label{subsec:two-moment-shield}
The linear contour perturbation in
Theorem~\ref{thm:rouche-shielding-independent}
can be improved without assuming a positive
moving zero gap. Let $t_*>\Lambda$, and
retain its three contour boundaries, common
minimum modulus $m>0$ and maximal imaginary
height $Y$. Define
\[
A_Y=\int_0^\infty\Phi(u)e^{t_*u^2+Yu}\,du,\qquad
M_Y=\int_0^\infty\Phi(u)e^{t_*u^2+Yu}u^2\,du.
\]
Both moments are strictly positive and finite.

\begin{theorem}[Sharp independent two-moment bound]
\label{thm:two-moment-shield}
For $0\le\tau\le\sigma$ the actual theta
heat flow satisfies, on every selected contour,
\begin{equation}\label{eq:two-moment-shield}
\begin{aligned}
|H_{t_*-\tau}(z)-H_{t_*}(z)|
&\le A_Y(1-e^{-\sigma M_Y/A_Y})\\
&\le\frac{2A_Y\sigma M_Y}{2A_Y+\sigma M_Y}\\
&\le\sigma M_Y.
\end{aligned}
\end{equation}
The first bound is optimal among positive
measures with mass $A_Y$ and second moment
$M_Y$; an atom at $u=\sqrt{M_Y/A_Y}$
achieves equality in Jensen's step.
\end{theorem}
\begin{proof}
Write $w_Y(u)=\Phi(u)e^{t_*u^2+Yu}$.
Positivity of the kernel and
$|\cos(zu)|\le e^{Yu}$ give
\[
|H_{t_*-\tau}-H_{t_*}|
\le\int_0^\infty w_Y(u)(1-e^{-\tau u^2})\,du.
\]
The integrand increases with $\tau$.
Jensen's inequality for the convex function
$v\mapsto e^{-\sigma v}$, applied to
$w_Y\,du/A_Y$, yields the first bound.
For $x\ge0$, the elementary inequality
$1-e^{-x}\le2x/(2+x)\le x$
follows from $e^{x/2}\ge1+x/2$
by differentiation. Substitute
$x=\sigma M_Y/A_Y$ to finish.
\end{proof}

\begin{corollary}[A longer noncircular
zero-pair continuation interval]
\label{cor:two-moment-window}
For $0<m<A_Y$ and $0<\sigma\le t_*$,
the rational sufficient condition
\begin{equation}\label{eq:two-moment-threshold}
\sigma<\frac{2mA_Y}{M_Y(2A_Y-m)}
\end{equation}
preserves one real simple zero in each
disk and no intervening zero. Consequently
$g(\tau)\ge g_*-2r>0$ and
\[
4\int_0^\tau\mathcal D(v)\,dv
\le g_*^2-(g_*-2r)^2<g_*^2.
\]
The exact exponential threshold is
$\sigma<(A_Y/M_Y)\log(A_Y/(A_Y-m))$.
\end{corollary}
\begin{proof}
Condition \eqref{eq:two-moment-threshold}
is equivalent to
$2A_Y\sigma M_Y/(2A_Y+\sigma M_Y)<m$.
Rouch\'e's theorem and real symmetry
then give the zero-count conclusions.
Only after establishing a gap floor do
we invoke $G'=-4\mathcal D$.
\end{proof}

For example, the rational moments
$A_Y=10$, $M_Y=20$, $m=6$, and
$\sigma=2/5$ fail the former linear
test ($\sigma M_Y=8>6$) but satisfy
the stronger one
($2A_Y\sigma M_Y/(2A_Y+\sigma M_Y)
=40/7<6$). A dependency-free
verifier checks 25,000 additional
exact rational gates. A
35-digit exploratory theta
calculation at $t_*=1/2$
gives $m_{\rm sample}/A_Y\approx
3.36\cdot10^{-4}$ for a
radius-$1/4$ circle about the first
positive zero. This is a
\emph{sampled upper bound} for the
continuous minimum, not a
certified contour margin;
the resulting nonlinear
improvement is small there.

The bound is strictly stronger as a
pair-dependent local theorem but does
not give boundary moduli for every
actual theta pair over the entire
backward heat interval. In particular
it does not exclude real-root loss
at unbounded height, or prove RH.
The algebraic implications are checked in
\path{XiKernel/TwoMomentShielding.lean}
and the rational verifier
\path{certificates/xi_two_moment_shielding_exact.py}.
Jensen's integral inequality and
the Rouch\'e transfer are manuscript
proofs, not end-to-end Lean theorems.
